\subsection{余代数的本原元素}

本号中我们考察一个余代数 E（《代数》，第 III 章，第 11 节，第 1 号），它带有余积

\[
c \colon \mathrm{E} \rightarrow \mathrm{E} \otimes \mathrm{E}
\]

和余单位 \(\varepsilon\)（同处，第 2 号）。回忆 \(\varepsilon\) 是 K-模 E 上的线性形式，满足（在 \(\mathrm{E} \otimes \mathrm{K}\) 与 \(\mathrm{K} \otimes \mathrm{E}\) 和 E 的典范等同下）：

\[
\mathrm{Id} _ {\mathrm{E}} = (\varepsilon \otimes \mathrm{Id} _ {\mathrm{E}}) \circ c = (\mathrm{Id} _ {\mathrm{E}} \otimes \varepsilon) \circ c.
\]

设 \(E^{+}\) 表示 \(\varepsilon\) 的核，并设 u 为 E 的元素，使得

\[
c (u) = u \otimes u \quad \text { and } \quad \varepsilon (u) = 1.
\]

K-模 E 是 \(E^{+}\) 与子模 K.u 的直和，后者是以 u 为基的自由模；设 \(\pi_{u}: E \to E^{+}\) 与 \(\eta_{u}: E \to K.u\) 为与这一分解相伴的射影。那么

\[
\pi_ {u} (x) = x - \varepsilon (x). u, \quad \eta_ {u} (x) = \varepsilon (x). u.\tag{1}
\]

定义 1. 元素 \(x\)（\(\mathbf{E}\) 的元素）称为 \(u\) -本原的，如果

\[
c (x) = x \otimes u + u \otimes x.\tag{2}
\]

\(u\) -本原的诸元素（\(\mathbf{E}\) 的）构成 \(\mathbf{E}\) 的一个子模，记作 \(\mathrm{P}_{u}(\mathrm{E})\)。

命题 1. 每个 \(u\) -本原元素（\(\mathbf{E}\) 的）都属于 \(\mathbf{E}^{+}\)。

\begin{enumerate}
\def\labelenumi{(\arabic{enumi})}
\setcounter{enumi}{1}
\tightlist
\item
  蕴含 \(x = \varepsilon(x).u + \varepsilon(u).x = \varepsilon(x).u + x\)，由此 \(\varepsilon(x) = 0\)。
\end{enumerate}

注. 若 \(x \in \mathrm{E}\) 且 \(c(x) = x' \otimes u + u \otimes x''\)，其中 \(x', x''\) 属于 \(\mathrm{E}^+\)，则 \(x = \varepsilon(x')\) . \(u + \varepsilon(u)\) . \(x'' = x''\)；同理 \(x = x'\)，于是 \(x\) 是 \(u\) -本原的。

对一切 \(x \in \mathrm{E}^{+}\)，我们记

\[
c _ {u} ^ {+} (x) = c (x) - x \otimes u - u \otimes x.\tag{3}
\]

命题 2. 我们有

\[
\left(\pi_ {u} \otimes \pi_ {u}\right) \circ c = c _ {u} ^ {+} \circ \pi_ {u}.\tag{4}
\]

设 \(x\) 属于 E；那么

\[
\begin{array}{r l} (\pi_ {u} \otimes \pi_ {u}) (c (x)) & = ((1 - \eta_ {u}) \otimes (1 - \eta_ {u})) (c (x)) \\ & = c (x) - (1 \otimes \eta_ {u}) (c (x)) - (\eta_ {u} \otimes 1) (c (x)) + (\eta_ {u} \otimes \eta_ {u}) (c (x)). \end{array}
\]

由于 \(\varepsilon\) 是 E 的余单位，

\[
(1 \otimes \eta_ {u}) (c (x)) = x \otimes u, \quad (\eta_ {u} \otimes 1) (c (x)) = u \otimes x
\]

由此

\[
\left(\eta_ {u} \otimes \eta_ {u}\right) (c (x)) = \left(\eta_ {u} \otimes 1\right) \left(\left(1 \otimes \eta_ {u}\right) (c (x))\right) = \varepsilon (x). u \otimes u;
\]

由此我们得出

\[
\left(\pi_ {u} \otimes \pi_ {u}\right) (c (x)) = c (x) - x \otimes u - u \otimes x + \varepsilon (x). u \otimes u.
\]

另一方面，

\[
c _ {u} ^ {+} \left(\pi_ {u} (x)\right) = c (x) - x \otimes u - u \otimes x + \varepsilon (x). u \otimes u,
\]

由此得公式 (4)。

由于 \(\mathrm{E}^{+}\) 是 \(\mathrm{E}\) 的直因子子模，\(\mathrm{E}^{+} \otimes \mathrm{E}^{+}\) 可等同于 \(\mathrm{E} \otimes \mathrm{E}\) 的一个直因子子模。在这一等同下，\(\pi_{u} \otimes \pi_{u}\) 是 \(\mathrm{E} \otimes \mathrm{E}\) 到 \(\mathrm{E}^{+} \otimes \mathrm{E}^{+}\) 上的射影。由公式 (4)，\(c_{u}^{+}\) 把 \(\mathrm{E}^{+}\) 映入 \(\mathrm{E}^{+} \otimes \mathrm{E}^{+}\)，且 \(\pi_{u}\) 是余代数 \((\mathrm{E}, c)\) 到余代数 \((\mathrm{E}^{+}, c_{u}^{+})\) 的态射。

命题 3. 若余代数 \((\mathrm{E}, c)\) 是余结合的（相应地，余交换的）（《代数》，第 III 章，第 11 节，第 2 号），则余代数 \((\mathrm{E}^{+}, c_{\mathbf{u}}^{+})\) 亦然。这由下述引理得出：

引理 1. 设 \(\pi: \mathrm{E} \to \mathrm{E}'\) 为满射余代数态射。若 \(\mathrm{E}\) 是余结合的（相应地，余交换的），则 \(\mathrm{E}'\) 亦然。

设 B 为结合 K-代数；映射 \(f \mapsto f \circ \pi\) 是 \(\mathrm{Hom}_{\mathbb{K}}(\mathrm{E}', \mathrm{B})\) 到 \(\mathrm{Hom}_{\mathbb{K}}(\mathrm{E}, \mathrm{B})\) 的单射代数同态。于是只需应用《代数》第 III 章第 11 节第 2 号的命题 1（相应地，命题 2）。

